% ============================================================
%  CONCLUSION
% ============================================================
\section{Conclusion}

\subsection{Limitations}

Several limitations bound the scope of our conclusions.
First, all experiments are restricted to 5{,}000 training iterations which is approximately one-tenth of
the compute used in the NanoChat leaderboard submission.
Techniques such as MTP and AttnRes, which require more data and steps to fully manifest their
representational advantages, may show different results under a longer training budget.
Second, hardware mismatch between experiment groups (4 $\times$ H100 vs.\
2 $\times$ H100) due to the limited availability of GPU resources makes it difficult to compare with the leaderboard.
Third, the hyperparameters of Engram (layer positions, memory size, gate initialisation) were not
extensively tuned for our scale, and a more careful configuration search might help mitigate some of
the observed degradation.
% Lastly, our analysis of FP8's beneficial effect is speculative; a proper ablation comparing equal
% throughput in FP16 vs.\ FP8 would be needed to isolate the regularisation effect.

\subsection{Future Work}

Several directions emerge naturally from our findings.
\begin{itemize}
  \item \textbf{Longer training}
  Re-running MTP and AttnRes with 50{,}000+ iterations would determine whether their
  representational advantages eventually outweigh the added complexity.
  \item \textbf{Engram modifications}
  A full-scale hyperparameter tuning of the engram layer either using an extensive grid search or more sophisticated methods such as Bayesian optimization can be done to improve model performance
  \item \textbf{RLHF fine-tuning}
  The best performing configuration (FP8 baseline) could be further improved through
  reinforcement learning from human feedback (RLHF) to enhance conversational quality and
  instruction following.
  \item \textbf{combining modifications}
  Different architectural configurations such as residual attention and engram layers can be applied together to evaluate the overall improvement.
\end{itemize}

\subsection{Broader Impact}

Our work demonstrates a reproducible proof-of-concept for conducting controlled experiments on GPT-2-scale chatbots on our limited computational resources.
Identifying which recent innovations transfer to small-scale language modelling and which do not can provide practical guidance for researchers and practitioners seeking to deploy capable autoregressive models in resource-constrained settings for offline use on-device mobile
inference, and applications where the cost of a frontier-scale model is unjustifiable.
